\section{Introduction}
\label{sec:intro}

% [TA] The mid-term introduction was "too short". Aim for 4-5 paragraphs that answer:
% what is the problem, why it matters, why tools, why small models, why these two domains,
% why two pressure levels, what we did and found.
\begin{outline}
\item \textbf{Problem and why it matters (ethics).} Sycophancy: a model drops a correct answer, or accepts a false statement, to agree with the user~\citep{sharma2024,perez2023,wei2023}. \textbf{[TA]} Frame the target broadly: an assistant should reject any false claim or assumption from the user, not only a wrong final answer. Harm: confirms false beliefs; larger in medicine and education, where users may act on the reply.
\item \textbf{Why tools (the main question).} \textbf{[TA]} Make the tool aspect the centre of the question sentence. Assistants are deployed with retrieval, search and code execution~\citep{lewis2020,schick2023,yao2023}; tools are often presented as a remedy for wrong answers. But if a tool-equipped model still gives in, the answer \emph{looks} evidence-backed and may be trusted more. Separate three things: access (tool offered), use (tool called), trust (answer follows the evidence).
\item \textbf{Why small language models.} \textbf{[TA]} (asked twice). Points to choose from (use your own reasons): most sycophancy studies test large proprietary assistants (SycEval: ChatGPT-4o, Claude-Sonnet, Gemini-1.5-Pro~\citep{fanous2025}; five large assistants in \citet{sharma2024}); 3--4B models are the ones people can run locally (privacy, cost, offline use in clinics or schools); small models follow tool-call formats less reliably~\citep{qwen2025,grattafiori2024,abouelenin2025}, so whether they choose and trust a tool is an open question; they fit the compute available to you.
\item \textbf{Why MedQA and MATH.} \textbf{[TA]} (asked in two comments). Points: both have one verifiable correct answer, so a change is measurable without a judge; they contrast two kinds of checking --- MATH claims can be checked by \emph{computation} (hence the Python tool), MedQA claims need \emph{knowledge} (hence retrieval, web, textbooks); medicine is the high-stakes case for an ethics course; SycEval used maths and medicine too, but with AMPS and free-text MedQuad --- multiple-choice MedQA lets you name the exact wrong option the user pushes. Say whether you expect results to differ from prior work and why (smaller models, multiple-choice format, tools).
\item \textbf{Why two pressure levels.} \textbf{[TA]} P1 is a plain suggestion (``I think the answer is \textit{X}, am I right?'', close to the user-belief prompts of \citet{sharma2024}); P2 adds claimed expertise (close to SycEval's ``ethos'' rebuttal~\citep{fanous2025}). You stopped at authority on purpose: SycEval's stronger levels add a fabricated justification or citation, i.e.\ \emph{fake evidence}, which would compete with the tool's evidence and blur what you measure (tool evidence vs.\ social pressure). Also each extra level multiplies runs (3 models $\times$ 6 conditions). Give your own reasons.
\item \textbf{What we did (2--3 sentences).} Define the conditions in words before their codes \textbf{[TA]}: no tool (T0), BM25 answer sheet with two entries (T1), web search (T2), the exact answer-sheet entry (Exact), textbook passages (T3), Python execution (Py); same 361 questions and the same wrong answer everywhere; paired tests.
\item \textbf{Findings (short list, with numbers).}
  (1) Without tools claimed expertise is the breaking point (MATH expert-turn resistance: Qwen \mathQwenPtwoTzero\%, Llama \mathLlamaPtwoTzero\%).
  (2) Tools help only when their output settles the disputed claim: Python raises expert-turn resistance for all models (+\deltaQwenPyPtwo, +\deltaLlamaPyPtwo, +\deltaPhiPyPtwo{} points); the exact entry helps Phi-4; web snippets and textbook passages can make things worse.
  (3) Access $\neq$ use $\neq$ trust: Qwen's tool choice is a position bias; Llama overrides correct evidence.
  (4) Implications for developers (Section~\ref{sec:implications}).
\end{outline}
